\chapterimage{numbers_background.pdf}

\chapter{Travessant el vel analògic-digital}\label{sec:piercing}

Enviem i rebem \conceptRef{message}{missatges} constantment en els nostres àmbits personal, professional i polític.
La disponibilitat gairebé instantània de pràcticament qualsevol \conceptRef{information}{informació} desitjada és una característica clau de l'era actual.
% 
Quan fem servir un ordinador de qualsevol mena per compartir o obtenir coneixement, primer cal transformar-lo en \conceptRef{data}{dades} que puguin entendre tant els humans com els ordinadors.
Això es pot imaginar com una seqüència de passos:
\begin{center}
\includegraphics[width=\linewidth]{knowledge_info_data_transmission.pdf}
\end{center}

Els ordinadors treballen amb \conceptRef{data}{dades} digitals i les transmeten, \ie, seqüències binàries
de \mbox{\zero\ i \one} com a símbols. 
Tanmateix, la transmissió de dades té lloc en el món real, on \mbox{\zero\ i \one} són idees i no objectes tangibles.

Aquest capítol estudia com aquests símbols digitals \textit{travessen el vel}  
del món digital al món analògic (físic) en la transmissió, i de tornada de l'analògic al digital en la recepció.


\section{Informació $\leftrightarrow$ Dades $\leftrightarrow$ Missatges binaris}\label{sec:piercing:information_binary}

El coneixement, la \conceptRef{information}{informació} i les \conceptRef{data}{dades} no són el mateix, 
tot i que sovint s'utilitzen indistintament. En contextos tècnics, cal utilitzar-los i interpretar-los amb precisió.
% 
\begin{exercise}
Suposem que som a l'aire lliure i volem explicar al nostre amic quin temps fa.
% 
Tenim coneixement de tot el que està relacionat amb el temps, perquè som allà.
Per a la nostra comunicació, ens hem de centrar en una part d'aquest coneixement i decidir quina informació enviem.
Per exemple, podríem voler comunicar la temperatura i dir alguna cosa com ``la temperatura és d'uns 22ºC''.
% 
Hi ha moltes maneres d'emmagatzemar aquesta informació com a dades dins d'un ordinador. El més probable és que representem el nombre 
$22$ com a part d'aquestes dades. Com definiries coneixement, informació i dades?
\end{exercise}


Tot i que hi ha molts tipus de dades, \ie\ numèriques, textuals, visuals, científiques, \etc, al final totes
es representen amb nombres. Per exemple, es podria fer servir la codificació \concept{ASCII} per representar la cadena ``Bobby''
com la seqüència següent: $66, 111, 98, 98, 121$.

Dins d'un ordinador, aquests nombres s'emmagatzemen en registres binaris, \eg\ a la CPU i a la RAM. 
Tanmateix, hi ha diverses maneres de representar dades numèriques en format binari. Més concretament,
hem de parar molta atenció, com a mínim, als aspectes següents sempre que llegim o escrivim dades:
\begin{itemize}
\item enter o amb decimals? 
\item \conceptRef{bitdepth}{profunditat de bits}? (\eg, 8 bits per nombre)
\item \conceptRef{sign}{amb signe} o \conceptRef{sign}{sense signe}?
\item Per a nombres de diversos bytes: \conceptRef{big endian/little endian}{big endian} o \conceptRef{big endian/little endian}{little endian}?
\item Per a nombres amb decimals: \conceptRef{fixed-point/floating-point}{coma flotant} o \conceptRef{fixed-point/floating-point}{coma fixa}?
\end{itemize}
% 
\begin{exercise}
Explica com completar la segona fila a partir de la primera (i viceversa),
i quines suposicions has de fer en cada cas.
% 
\begin{center}
\begin{tabular}{lcccccc}
\toprule
\textbf{Decimal}
& $7$
& $2$
& $-3$
& $0.75$
& $260$
\\
\midrule
\textbf{Binari}
& \mbox{\zero\zero\one\one\one}
& \mbox{\zero\zero\zero\zero\zero\zero\one\zero}
& \mbox{\one\one\zero\one}
& \mbox{\one\one\zero\zero}
& \mbox{\zero\zero\zero\zero\ \zero\one\zero\zero\ \zero\zero\zero\zero\ \zero\zero\zero\one}
\\
\bottomrule
\end{tabular}
\end{center}
% 
I què passa amb à $\leftrightarrow$ \one\one\zero\zero\ \zero\zero\one\one\ \one\zero\one\zero\ \zero\zero\zero\zero?
\end{exercise}


De vegades hem d'inspeccionar el contingut d'un registre binari 
(\eg, per depurar o amb altres finalitats d'anàlisi) o establir-ne el contingut manualment.
% 
Per als humans, és incòmode i \textit{molt} propens a errors 
manipular cadenes binàries de més d'uns quants bits. 
Quan hem d'inspeccionar o modificar el contingut d'un registre binari
(\eg, per depurar), sovint fem servir la base $16$, \ie, la notació \concept{hexadecimal}.


En hexadecimal, hi ha exactament $16 = 2^4$ dígits diferents (de \otherBase{0} a \otherBase{9}, de \otherBase{a} a \otherBase{f})
amb valors decimals de $0$ a $15$, cadascun dels quals representa exactament $4$ bits.
Quan escrivim valors hexadecimals per a un ordinador, \ie, en codi font, els valors hexadecimals
normalment van precedits de \otherBase{0x}, \eg, \otherBase{0x58a5b0}. 
De manera semblant, les expressions binàries van precedides de \otherBase{0b}, \eg, \otherBase{0b0011}.

\begin{remark}
\begin{minipage}{\linewidth}
\begin{itemize}[leftmargin=0.5cm]
\item Els espais i fins i tot els salts de línia entre dígits no tenen cap significat; només s'utilitzen 
per facilitar la inspecció visual.
% 
\item En alguns textos, quan en un context hi ha diverses bases possibles, s'indiquen
com a subíndexs, com a \otherBase{58a5b0$_\textrm{16}$} i \otherBase{0011$_\textrm{2}$}.
\end{itemize}
\end{minipage}
\end{remark}

\begin{exercise}
Considera el missatge següent, format per una concatenació 
d'enters sense signe de 3 bits: \otherBase{a4 3f 20}.
\begin{itemize}
\item Quants bits i quants bytes fa?
\item Quins són els cinc primers enters continguts en el missatge?
\end{itemize}
\end{exercise}

La majoria de vegades, deixarem que el nostre codi s'encarregui de manipular les dades. Per fer-ho, necessitem
operacions \conceptRef{bitwise}{a nivell de bits} i de \conceptRef{boolean logic}{lògica booleana} com les següents

\begin{center}
\begin{tabular}{rcccccc}
\toprule
\textbf{Operació} & 
AND (bits) &
AND lògica &
OR (bits) &
OR lògica &
Despl. esq. &
Despl. dre.
\\
\textbf{Python}    &
\inlineCode{&}       & 
\inlineCode{and}    & 
\inlineCode{|}       & 
\inlineCode{or}    & 
\inlineCode{<<} &  
\inlineCode{>>} 
\\
\bottomrule
\end{tabular}
\end{center}

Altres eines útils de Python són les funcions \texttt{bin()} i \texttt{hex()}, que converteixen un enter 
en la seva representació binària i hexadecimal, respectivament, i \texttt{int()}, que pot analitzar
una cadena que descriu un nombre i retornar aquest nombre. També podem controlar el format amb què mostrem els nombres,
\eg, \inlineCode{print(f'{n:08b}')} mostraria el valor de la variable \inlineCode{n} en forma binària, 
fent servir com a mínim 8 posicions i omplint amb zeros les posicions buides de l'esquerra
(més informació a la \href{https://docs.python.org/3/reference/lexical_analysis.html#f-strings}{documentació de Python}).

\begin{exercise}
Considera el codi que es mostra a continuació. Els nombres que s'imprimeixen en executar-lo tenen una importància especial en xarxes i en informàtica.
Què tenen en comú? (Pista: potser et convé modificar el codi perquè en mostri les representacions binàries).

\begin{center}
\showCode{snippets/bitwisemanipulation.py}
\end{center}
\label{ex:bitwisemanipulation}
\end{exercise}

\section{Analògic $\leftrightarrow$ Digital}\label{sec:piercing:analog_digital}

Un cop decidim quins bits volem transmetre, hem de trobar una manera de fer que aquests bits arribin a una altra màquina a distància.
Primer, hem de decidir quin \conceptRef{medium}{mitjà} fem servir i quines propietats físiques podem controlar i monitorar.
Entre les opcions més habituals hi ha fer servir el voltatge en cables de coure, enviar llum per fibra òptica i manipular el 
camp electromagnètic mitjançant antenes.

\begin{exercise} Reflexiona:\\[-0.5cm]
\begin{itemize}
\item Podem fer un cable de coure tan llarg com vulguem?
\item Quin mitjà fan servir les connexions Wi-Fi? 
\item Necessitem fibres òptiques per fer comunicacions basades en la llum?
\end{itemize}
\end{exercise}


Si optem per cables de coure, en podem posar diversos en \conceptRef{parallel/serial}{paral·lel} i enviar més dades a la mateixa velocitat de rellotge.
Això, però, té el preu d'una complexitat i un cost més grans que els dissenys en \conceptRef{parallel/serial}{sèrie}. 
Es poden produir interferències de voltatge a les línies de dades per diversos motius, 
incloent-hi fenòmens físics com la inducció electromagnètica al cable. 
Tot i així, els cables que compleixen els estàndards moderns de transmissió de dades (\eg, IEEE 802.3)
sovint fan servir parells trenats de cables i apantallament,
entre altres estratègies, per garantir errors de bit per sota d'$1$~de~cada~$10^{12}$ bits.

\begin{exercise}\label{ex:piercing:copper_connectors}
Observa i identifica els connectors de cable següents:\\[-0.5cm]
\begin{center}
\includegraphics[height=3cm]{hdmi_connector.jpg}
\includegraphics[height=3cm]{coaxial_connector.jpg}
\includegraphics[height=3cm]{usbc_connector.jpg}
\includegraphics[height=3cm]{rj11_connector.jpg}
\end{center}
\end{exercise}

\begin{wrapfigure}{L}{0.22\linewidth}

\includegraphics[width=\linewidth]{iceberg.jpg}
\end{wrapfigure}
Alguna vegada t'has capbussat en una piscina, has mirat cap amunt i l'aigua semblava un mirall?
Els cables de fibra òptica funcionen gràcies a aquest fenomen, anomenat refracció.
La llum que entra a la fibra per un extrem 
s'hi queda a dins fins que surt per l'altre extrem. El recobriment de la fibra és important
per preservar-ne la integritat, però aquest recobriment no participa en el ``rebot'' de 
la llum mentre avança per la fibra.

En una fibra òptica monomode, el sensor de l'extrem receptor 
detecta la presència o l'absència d'un únic rang de longituds d'ona. Aquests cables són relativament prims i barats,
però transporten un sol senyal.
% 
Una altra opció és transmetre diverses longituds d'ona (\eg, fent servir diverses fonts làser) alhora, 
per la mateixa fibra. A l'extrem receptor, es fa servir un prisma per tornar a dividir el feix de llum en els seus
components monocromàtics, que es detecten per separat. 
D'aquesta manera, la fibra òptica multimode permet \conceptRef{multiplex}{multiplexar}
diversos senyals, cosa que augmenta molt la taxa de transmissió efectiva.


\begin{exercise}\label{ex:piercing:laser_refraction}
Quan un feix de llum entra en una fibra òptica, ho fa amb un cert angle d'incidència.
Aquest angle afecta el temps que triga a arribar a l'altre extrem. Si projectem un con de llum,
la llum entra a la fibra amb diversos angles d'incidència. 
Què passa a l'altre extrem si encenem i apaguem aquest con de llum molt ràpidament?
És rellevant si fem servir feixos monomode o multimode?
% 
\begin{center}
\includegraphics[height=2.5cm]{laser_refraction.jpg}
\includegraphics[height=2.5cm]{optical_connector.jpg}
\end{center}
\end{exercise}
\newpage

\begin{wrapfigure}{R}{0.3\linewidth}
\begin{center}
\includegraphics[width=\linewidth]{dipole.png}
\end{center}
\end{wrapfigure}

Tot i que hi ha moltes classes d'antenes, la seva finalitat principal 
és detectar camps electromagnètics i/o crear-ne. L'equació de Maxwell-Faraday ens diu
que els \textit{canvis} en el \conceptRef{field (physics)}{camp} elèctric $\vec{E}$ creen un \conceptRef{field (physics)}{camp} magnètic perpendicular $\vec{B}$,
i a l'inrevés.

Una antena de dipol com la de la dreta aprofita aquest concepte aplicant un voltatge altern $V$,
que genera ones electromagnètiques que es propaguen cap a fora de l'eix del dipol. 

\begin{remark}
Quan emetem una certa energia $E$ des d'un punt en totes direccions, aquesta energia es distribueix al voltant d'una esfera creixent.
Com que l'àrea d'una esfera de radi $r$ és $4 \pi r^2$, la quantitat d'energia que arriba a un punt a distància $r$
serà de l'ordre de magnitud de $E / r^2$.

Això de vegades s'anomena llei de l'invers del quadrat, i és
el motiu pel qual els corrents induïts a l'antena del receptor solen estar entre els nanovolts 
($10^{-9}$~V) i els picovolts ($10^{-12}$~V). Sorprenentment, això és suficient perquè els detectors moderns llegeixin el senyal de dades.
\end{remark}


\begin{wrapfigure}{R}{0.2\linewidth}
 \includegraphics[width=\linewidth]{modulation.jpg}
\end{wrapfigure}


Independentment del mètode triat perquè les dades viatgin, cal travessar igualment el vel entre els senyals analògics i les dades 
digitals. Una manera de fer-ho és la modulació: es genera un senyal sinusoïdal de base anomenat portadora,
i les dades es codifiquen modificant aquesta portadora. Per exemple, es pot canviar l'amplitud de la portadora (ASK),
la seva freqüència (FSK), o fins i tot l'amplitud i la fase alhora (QAM). 

\begin{exercise}
La figura anterior mostra un exemple de modulació d'amplitud i un de modulació de freqüència. 
Quin és quin? Per a cada cas: aquesta modulació transmet un senyal analògic o digital?
\end{exercise}


\begin{exercise}
Quina és la diferència entre amplada de banda i velocitat de transmissió? 
Per què creus que s'utilitzen indistintament en contextos no tècnics?
\end{exercise}
